\subsection{An actual finite congruence obstruction for coefficient erasures}

\begin{statement}{A.67}{The j-line levels three and four detect explicit rank-one ghosts}{Proved}
Use the normalized modular coordinate $J=j/1728$ on
$T=\mathbf P^1-\{0,1,\infty\}$, so that $X(m)\to\mathbf P^1_J$
has branch orders $(3,2,m)$. Let
\[
 E:\ \alpha X^3+\beta Y^3=Z^3,\qquad
 J=\alpha(X/Z)^3,\quad1-J=\beta(Y/Z)^3,
 \tag{A67.1}
\]
be a smooth diagonal cubic over $\Q$ with a rational flex
$O$ in the $Z=0$ block, chosen as origin. Put
$D=E[3]=\{XYZ=0\}$, $U=E-D$, and let $R\in E(\Q)$
be nontorsion. Write $\Pi_m$ for the arithmetic quotient
of $\pi_1(T)$ attached to the finite projective level-$m$
cover, with geometric group
$G_m=\mathrm{SL}_2(\Z/m\Z)/\{\pm I\}$. Define
\[
 \Gamma_m=
 \operatorname{im}\!\left[
  \pi_1(U)\longrightarrow
  \pi_1(E)\times_{G_\Q}\Pi_m\right].
 \tag{A67.2}
\]
This is the actual joint quotient retaining the entire
compact fundamental group and the indicated finite
modular quotient. A Selmer lift to $\Gamma_m$ means
a section projecting to $s_a$ on $\pi_1(E)$ and locally
equal, at every finite place, to the image of an
interior point or cusp section of $U$.

For every nonzero integer $n$ and nonempty finite set
$\mathcal P$ of coefficient primes, the following
hold.
\begin{enumerate}
\item If $3\notin\mathcal P$, the exponent
      $a=n(1-e_{\mathcal P})$ has no Selmer lift
      to $\Gamma_3$.
\item Without a restriction on $\mathcal P$,
      it has no Selmer lift to $\Gamma_4$.
      This obstruction already factors through
      the quadratic character $\kappa_2(1-J)$.
\item On the least-conductor rank-one cubic
      $E:y^2+y=x^3-1$, with $P=(1,0)$ and its
      infinity flex as origin, every singleton
      erasure $a=n(1-e_p)$, for every prime $p$
      and every $n\ne0$, has no Selmer lift to
      $\Gamma_3$.
\end{enumerate}
The first two assertions apply uniformly to
the elliptic cubic twists $X^3+Y^3=dZ^3$
with their rational infinity flex.
They do not silently identify a genus-one
torsor without a rational origin with its
Jacobian.

These are failures of Selmer lifting. All
the exponents still have abstract lifts
to the same finite-kernel quotients by A.44.
No complete lifting set for $\Gamma_3$ or
$\Gamma_4$, or for the full punctured modular
cover groups, is asserted.
\end{statement}
\begin{proof}
\par\smallskip\noindent\textbf{The exact geometric comparison.}\enspace 
Let $F=\pi_1(T_{\overline\Q})$ have peripheral
generators $A,B$ around $J=0,1$.
The Fermat subgroup is
$H=\ker(F\to(\Z/3\Z)^2)$, with the two
generators mapped independently.
For projective modular monodromy take
\[
 S=\begin{pmatrix}0&-1\\1&0\end{pmatrix},
 \quad T_0=\begin{pmatrix}1&1\\0&1\end{pmatrix},
 \qquad A=ST_0,\quad B=S,\quad C=(AB)^{-1}.
 \tag{A67.3}
\]
Modulo the centre their orders are $3,2,m$,
respectively. The image of $H$ is generated
by $G_m'$ and the cubes of $A,B$.
For $m=3$ it is $V_4=A_4'$; for $m=4$
it is all of $S_4$, since $B$ is nontrivial
in $S_4/A_4$; and for $m=5,7$ it is the
whole perfect group. The images of the
cubic cusp inertia groups are conjugates
of $A^3,B^3,C^3$. Their normal closure is
the entire indicated image: in $V_4$ the
three order-two conjugates generate, in
$S_4$ the transpositions generate, and
the two remaining groups are simple.

Thus the cusp kernel already surjects to
the finite geometric factor. The joint
geometric image in (A67.2) is consequently
the product of $T E$ with that factor;
in particular the kernel over $\pi_1(E)$
is finite. The connected compact fibre
products have the following data:
\[
\begin{array}{c|r|r|r|c|r}
 m&|G_m|&|\operatorname{im}H|&
 \text{components}&
 \text{inertia over }0,1,\infty&
 \text{genus per component}\\ \hline
 3&12&4&3&(1,2,1)&4\\
 4&24&24&1&(1,2,4)&46\\
 5&60&60&1&(1,2,5)&118\\
 7&168&168&1&(1,2,7)&343
\end{array}
\tag{A67.4}
\]
There are three cubic cusps in each
branch block; Riemann--Hurwitz gives
$2g-2=\sum 3|\operatorname{im}H|(1-1/e)$.
These are genera of the fibre products
with the compact cubic, not the genera
of $X(m)$ themselves.

For clarity, the actual genus-zero
$X(3)$ map is explicit. The Hesse family
$x^3+y^3+z^3-3kxyz=0$, with its constant
nine flexes, has Weierstrass model
\[
 Y^2=X^3-27k(k^3+8)X+
                54(k^6-20k^3-8).
\]
Writing $W=x+y+kz$, one transformation
is $X=12(k^3-1)z/W-9k^2$ and
$Y=36(k^3-1)(x-y)/W$. Its invariant gives
\[
 J=\frac{k^3(k^3+8)^3}{64(k^3-1)^3},
 \qquad
 1-J=-\frac{(k^6-20k^3-8)^2}
                       {64(k^3-1)^3}.
 \tag{A67.5}
\]
The degree is twelve, with ramification
partitions $3^4,2^6,3^4$, as required
for the full projective level-three
cover. On the untwisted cubic a connected
pullback component is
\[
 h^3=k^6-20k^3-8,\quad
 u=\frac{k(k^3+8)}{4(k^3-1)},\quad
 v=-\frac{h^2}{4(k^3-1)},\quad u^3+v^3=1.
 \tag{A67.6}
\]
For (A67.1), after choosing
$\alpha^{1/3}$, replace this by
$h^3=\beta(k^6-20k^3-8)$, divide $u$
by $\alpha^{1/3}$, and divide $v$ by
$\beta$. The sextic polynomial is
squarefree; its degree-three cover
has genus four, also directly verifying
the first row of (A67.4).

\par\smallskip\noindent\textbf{The quadratic characters actually retained.}\enspace 
Let $D_1$ be the three cusps above
$J=1$. In the $V_4$ cover over the
cubic, their inertia elements are the
three distinct nonzero elements.
Each nontrivial $V_4$ character therefore
ramifies at exactly two of them. Its
intermediate double cover has genus two.
Over a finite extension $L$ splitting
$E[3]$, for any distinct $d_i,d_j\in D_1$
choose a unit
\[
 f_{ij}\in L(U)^\times,\qquad
       \operatorname{div}(f_{ij})
                     =3(d_i)-3(d_j).
 \tag{A67.7}
\]
Its mod-two geometric cusp character is
the corresponding $V_4$ character.
The difference of these two geometric
characters kills all cusp inertia and
therefore factors through
$\pi_1(E_{\overline L})$.

Here is the arithmetic descent argument,
without assuming an unwarranted extension
of an invariant geometric character.
The Kummer character $\kappa_2(f_{ij})$
is already an arithmetic character on
$\pi_1(U_L)$. The kernel of
$\pi_1(U_L)\to\Gamma_{3,L}$ is geometric
and is killed by both compactification
and the modular quotient. The preceding
geometric comparison shows that
$\kappa_2(f_{ij})$ kills this kernel.
It therefore factors through
$\Gamma_{3,L}$ itself. A global Selmer
section supplies a single global
mod-two class with local evaluations
$f_{ij}(Q_v)$ at every interior centre.
No larger extension than a field
defining the functions is needed.

For level four, the quotient
$S_4\to\Z/2$ is unramified at $J=0$
and ramified at $J=1,\infty$.
Its character is
$\kappa_2(1-J)$ up to a global
quadratic constant. Thus a
$\Gamma_4$ section supplies the
global mod-two class of that unit.
On the cubic,
\[
 \operatorname{div}(1-J)
       =3\sum_{d\in D_1}(d)
          -3\sum_{d\in D_\infty}(d).
 \tag{A67.8}
\]
In particular its pole at the chosen
origin $O\in D_\infty$ has odd order.

\par\smallskip\noindent\textbf{Uniform erasures.}\enspace 
Set $R'=nR$, $M=\prod_{p\in\mathcal P}p$.
At a finite place $v$ away from
$\mathcal P$ write
\[
 R'=Q_v+T_v,\qquad
 Q_v=(1-e_{\mathcal P})R',\quad
 T_v=e_{\mathcal P}R'.
 \tag{A67.9}
\]
Choose one finite set $S$ containing
$3$, $\mathcal P$, bad reduction,
the ramification of $L$, vertical
divisor components, and representatives
of the finitely many global quadratic
classes just obtained. Outside $S$,
$T_v$ is finite torsion of order
prime to the residue characteristic,
whereas $Q_v$ is nontorsion. Full
compact local Kummer injectivity
forces $Q_v$ to be the actual
interior centre of a Selmer lift.
Consequently the indicated unit
evaluations have even valuation
outside $S$.

Consider the infinitely many distinct
points $R_j=[M^{2j}]R'$. If one reduces
to $d\in E[3]$, primary decomposition,
followed by injectivity of reduction
on prime-to-residue-characteristic
torsion, gives $d_0,d_1\in E[3](L)$
with
\[
 \overline{Q_v}=\overline{d_0},\quad
 [M^{2j}]T_v=d_1,\quad
 d=[M^{2j}]d_0+d_1,\quad
 R_j-d=[M^{2j}](Q_v-d_0).
 \tag{A67.10}
\]
The two positive formal distances
in the last equality are equal,
since $M$ is a local unit.

If $3\notin\mathcal P$, then $d_1=O$:
it is simultaneously
$\mathcal P$-primary and $3$-torsion.
Also $M^{2j}\equiv1\pmod3$, so
$d=d_0$. Whenever $R_j$ approaches
a cusp in $D_1$, $Q_v$ approaches
the same cusp with the same distance.
At least one character (A67.7)
has odd order there, so the even
valuation of its local evaluation
makes that distance even.
Fix one $f_{ij}$. Its values on
every $R_j$ have even valuation
outside $S$. They lie in finitely
many $S$-squareclasses. The
corresponding twists of
$z^2=f_{ij}$ have genus two, so
Faltings contradicts the infinitely
many $R_j$
\citep[Satz~7, p.~365, with the correction]{AMFaltings1983,AMFaltings1984Erratum}.
This proves the first assertion.

For $\Gamma_4$ use $f=1-J$ alone.
If $3\notin\mathcal P$, the same
argument applies whenever $R_j$
approaches its zero or pole divisor.
If $3\in\mathcal P$, the retained
component $Q_v$ has no $3$-primary
reduction. Hence the $3$-torsion
point $d_0$ in (A67.10) is $O$.
The function $f$ has an odd pole
there, so its even valuation at
$Q_v$ makes the distance in
(A67.10) even. In either case
$v(f(R_j))$ is zero or plus or
minus three times an even
distance, outside the same $S$.
The finitely many twists of
$z^2=f$ now have six branch
points and genus four.
Faltings gives the same
contradiction, proving the
second assertion. This explains
exactly the origin condition:
the proof uses the odd pole at
the rational infinity flex.

\par\smallskip\noindent\textbf{The erased prime three on the least cubic.}\enspace 
It remains to prove the third
assertion for $p=3$; the other
primes follow from the first.
On $E:y^2+y=x^3-1$ the Fermat
coordinates are
\[
 u=(y+2)/x,\quad v=(1-y)/x,\quad
 u^3+v^3=9,\quad J=u^3/9.
 \tag{A67.11}
\]
Put $b^3=3$, $\zeta^2+\zeta+1=0$,
and $L=\Q(b,\zeta)=\Q(E[3])$.
The three cusps of $D_1$ are
$(b\zeta^i,1)$, and a character
of the form (A67.7) is
\[
 g=\frac{b^2\zeta^2x-y-2}{b^2x-y-2}.
 \tag{A67.12}
\]
We first establish a forcing fact.
If a $\Gamma_{3,L}$ Selmer
section has compact class
$(1-e_3)\kappa(A)$ for a
nontorsion $A\in E(L)$, then
its global $g$-class is
$\kappa_2(g(A))$.
The complete $3$-torsion and
$A$-Kummer tower over $L$ is
pro-$3$: the residual torsion
action is trivial, the higher
torsion action lies in the
pro-$3$ congruence kernel,
and the translation module
is $T_3E$. A.61 gives a
nonempty open Frobenius
condition \emph{inside this
tower over $L$} where the
local $3$-primary component
of $A$ is zero.

A nonzero quadratic difference
between the proposed global
$g$-class and $\kappa_2(g(A))$
would define a quadratic
extension disjoint from
this pro-$3$ tower.
Chebotarev would give a good
place of residue characteristic
different from three, outside
all finite supports, with
nontrivial quadratic Frobenius
and with that zero-primary
condition. At this place
$(1-e_3)A=A$ as an actual
local point: the local
$3$-primary Kummer class
vanishes and reduction is
injective on such torsion.
The local unit difference
is therefore trivial, a
contradiction. This proves
the forcing fact.
The unramified compact-character
comparison introduces no loss;
alternatively its values agree
on these compact classes because
their coefficient-two components
agree. The direct kernel
argument preceding (A67.8)
already proves that $g$ itself
is retained by the joint quotient.

Now let $d_0=(b,1)$,
$d_1=(b\zeta,1)$. One has
$d_0-d_1=(0,\zeta)$. Translation
in (A67.12), followed by a
divisor comparison, gives the
exact identity
\[
 g(A+d_1)=\frac{\zeta+2}{\zeta-y(A)}.
 \tag{A67.13}
\]
Both sides have divisor
$3(O)-3((0,\zeta))$ as
functions of $A$. Evaluating
at $A=-d_1$, using $g(O)=1$
and $y(-d_1)=-2$, determines
the constant. Thus no
asymptotic approximation is
involved.

There is a unique place $w$
of $L$ over three, with
$e(L_w/\Q_3)=6$ and
$v_w(\zeta+2)=3$. For a finite
point $A=(x,y)\in E(\Q_3)$,
integral $y$ cannot satisfy
$y\equiv1\pmod3$: then
$y^2+y+1$ has $3$-adic
valuation one, whereas
$x^3$ has valuation divisible
by three. For the other
integral residues,
$v_w(\zeta-y)=0$.
If $y$ is nonintegral,
$v_w(\zeta-y)=6v_3(y)$.
Thus this valuation is
always even. Moreover
$v_w(g(A))=0$, since its
numerator and denominator
are conjugate under the
local automorphism
$b\mapsto b\zeta$, which
fixes the rational coordinates.
Equation (A67.13) proves
\[
 v_w\!\left(\frac{g(A+d_1)}{g(A)}\right)
       =3-v_w(\zeta-y(A))\quad
                   \text{is odd}.
 \tag{A67.14}
\]
The ratio is not a local square.

Take $A=nP$, with $n\ne0$.
Since $(1-e_3)$ kills the
full Kummer class of the
$3$-torsion point $d_1$,
the compact classes
$(1-e_3)\kappa(A)$ and
$(1-e_3)\kappa(A+d_1)$
are identical over $L$.
The forcing fact applies
to both nontorsion reference
points. A Selmer lift would
make its single global
$g$-class equal both
$\kappa_2(g(A))$ and
$\kappa_2(g(A+d_1))$.
This contradicts (A67.14),
and proves the assertion
for every nonzero $n$.
\end{proof}

There is a small independent finite
witness for the last argument at
$n=1$. At $q=61$ choose the
embedding $L\to\Q_{61}$ with
$b\equiv4$, $\zeta\equiv13$.
Exact addition on $E(\mathbf F_{61})$
gives
\[
\begin{array}{c|c}
 \operatorname{ord}(\overline P)&21\\
 \text{$3$-primary component of }\overline P&(5,59)\\
 \overline{(1-e_3)P}&(36,36)\\
 g(\overline P),\ g(\overline{(1-e_3)P})&10,\ 60\\
 g(\overline{(1-e_3)P})/g(\overline P)&6,\quad6^{30}=-1.
\end{array}
\tag{A67.15}
\]
Thus the forced global equality also
fails by a nonsquare residue at a
good split prime. The exact scripts
\texttt{a\_main\_jline\_covers.py},
\texttt{a\_main\_jline\_at3.py} and
\texttt{a\_main\_jline\_local3.gp}
verify the modular matrix groups,
the Hesse and translated-function
identities, the genera, this finite
witness, and local valuation samples.
Their complete executed outputs
are preserved; the uniform proof
does not extrapolate from the samples.

The mechanism at level four is
already the degree-two cover
$w^2=1-J$. It is also the sign
quotient of the $S_3$ cover at
modular level two. Thus neither
the full order-twenty-four group
nor an unspecified interaction
of coefficient primes is
essential. Among the tests named
in task A.12, level four gives
the uniform bound above, while
level three suffices for every
singleton erasure on the least
curve. These modular levels
are not nested unless one takes
their joint refinements; no
minimality claim is inferred.
Failure in one of the explicitly
identified joint quotients
excludes a compatible lift to
a larger group retaining it.
The converse does not hold.
The following A.68 extends the
level-three exclusion to every
finite alteration of a nonzero
integer on the least curve.
General profinite exponents
and full metabelian and
punctured-cover fibres remain
outside these results.
